% !TEX root = ../../main.tex
% Part II (Nguyen Huu Cuong) - merged from the research report (report_latex_v3), see MERGE_NOTES.md.
% Style pass 2026-09-30: em dashes and rhetorical sentences removed; technical content, numbers and references unchanged.
\chapter{Experiments and Results}
\label{chapter:r-experiments}

\textit{This chapter presents the experimental setup on Market-1501, the results of the reasoning layer against the CLIP baseline with an analysis of what it rescues and what it loses, the ablations of its mechanisms, and the tooling that makes every number reproducible.}

\section{Experimental Setup}

\subsection{Dataset}
\label{sec:r-dataset}
The experiments use Market-1501~\cite{zheng2015market1501}, with the attribute
labels of the Market-1501\_\allowbreak Attribute project~\cite{lin2019market1501attr}. The
dataset contains 1501 identities, captured by six cameras on a university campus,
and its person boxes were produced by an automatic pedestrian detector rather than
drawn by hand.

\paragraph{Why Market-1501, and how it stands in for a tracker.} The reasoning
layer needs several crops of the same person, grouped together, so that the frames
of a track can agree or disagree about an attribute. Market-1501 provides this,
because every crop carries an identity label. Grouping crops by identity gives each
person a set of observations under different viewpoints, lighting, blur and
occlusion, which is what a tracker passes to the search block in the deployed
pipeline.

This is not a single-camera setting. A single-camera tracker links a person across
consecutive frames of one view, whereas the identity groups of Market-1501 span up
to six cameras. They therefore play the role of an \emph{ideal} tracker, one that
never switches identities, never merges two people, and follows a person across
cameras. The views inside one group differ more than the frames of a single-camera
track would, so disagreement between readings, which the funnel is built to reason
about, is more frequent rather than less. Market-1501 also carries third-party
hand annotation of attributes~\cite{lin2019market1501attr}, which makes an
attribute-level error analysis possible.

\paragraph{Why YOLO11n and ByteTrack are not in the experiment.} The experiment
tests the search block in isolation: the encoder, the retrieval of the candidate
pool, and the reasoning funnel behind it. Detection and tracking are left out on
purpose. Market-1501 already supplies the detected person crops and the identity
grouping, i.e.\ the outputs those two blocks would produce, so running YOLO11n and
ByteTrack on top would add nothing the dataset does not already provide, while
mixing their errors into the measurement. With them removed, every failure
reported in Section~\ref{sec:r-error} belongs to retrieval or to reasoning. The
cost of this isolation is that upstream errors are not measured at all;
Section~\ref{sec:r-limits} returns to this point. YOLO11n and ByteTrack remain the
chosen blocks of the deployed pipeline (Table~\ref{tab:r-choices}).

\paragraph{The 496-identity subset.} From the dataset and its hand annotation we
keep only the identities that are \emph{uniquely identified} by their annotated
attributes, that is, no other identity carries the same attribute vector. This
reduces 1501 identities to 496, and every rate in this part is drawn from that
population.

The reason is that the task has to be answerable from the description alone. If
two people in the gallery are described identically, a text query giving that
description has two correct answers while the annotation records only one, so a
system that returns the other is marked wrong although it is right. Removing those
identities removes a source of noise that the system cannot act on. It also makes
the benchmark \emph{easier} than the full 1501 identities, which must be taken into
account whenever these numbers are compared with published Market-1501 results,
since those are not computed on this subset.

\subsection{Models}
The models used in the experiment are the search-block entries of
Table~\ref{tab:r-choices}: CLIP ViT-H/14 (laion2B) as the base encoder, Milvus as
the vector database, Qwen2.5-VL-3B-Instruct as the attribute describer, and Gemini
Flash-Lite 3.5 as the reasoning LLM at thinking level HIGH for every stage.
Detection and tracking are supplied by the dataset itself
(Section~\ref{sec:r-dataset}), so YOLO11n and ByteTrack, the deployed pipeline's
choices for those blocks, are not run here.

No model is fine-tuned on this dataset, and every number below is zero-shot.
Sampling temperature is left at the provider default, following the guidance for
this model family.

\subsection{Funnel sizes and token budget}
To fit the token limits of the free Google AI Studio tier, we retrieve a cached
pool of 150 crops and use the top 75; the coarse round keeps 25 and the final
round returns 10. Stage 2A is split into three calls of 25 candidates.

\begin{table}[H]
\centering
\begin{tabular}{@{}lrr@{}}
\toprule
\textbf{Stage} & \textbf{Input tokens (est.)} & \textbf{Output tokens}\\
\midrule
1 and 3                  & $\sim$3k          & $\sim$200\\
2A (per split) and 4A    & 40k $\rightarrow$ 60k & 15k $\rightarrow$ 20k\\
2B                       & $\sim$60k         & $\sim$2k\\
4B                       & $\sim$20k         & $\sim$1k\\
\bottomrule
\end{tabular}
\caption{Per-call token budget. Output excludes thinking tokens, which are billed
as output and counted against the same rate window, so the real output load is
higher than this column. The reliable signal for a truncated reply is the API's
\code{finish\_reason}, not the reported output count.}
\label{tab:r-tokens}
\end{table}

\subsection{Baseline}
\paragraph{image-clip.} The baseline returns the top-10 \emph{crops} of the
75-candidate pool by cosine similarity, ungrouped, with repeats kept. Every table
reports this baseline. It is a like-for-like comparison: it returns the same kind
of result as the funnel, from the same pool.

\paragraph{A grouped baseline, considered and not reported.} An earlier sweep
also scored a \emph{track-naive} row, which assumes a perfect tracker: the pool's
crops are grouped by track, each person is scored by the maximum cosine over their
crops blended with how many of their crops surfaced
($0.7\,s_{\text{norm}} + 0.3\,c_{\text{norm}}$, both normalised within the
query's own pool), and the people are ranked by this score. It is not carried into
the results below, for two reasons. It answers a different question: it returns
ten distinct \emph{people} where image-clip and the funnel return ten
\emph{crops}, so its R@10 has a different denominator of opportunity. In addition,
on the earlier 44-identity sweep where it was measured, it tied image-clip exactly
at every $k$ and recovered no identity that CLIP's raw ordering had missed. This
cheaper grouping alternative to the reasoning layer therefore did not help.

\subsection{Metrics}
\label{sec:r-metrics}
R@$k$ is the fraction of queries whose ground-truth identity appears in the
returned top-$k$. Alongside it we report \code{cap@75}, how often the ground truth
was in the 75-crop pool at all. No reranking can exceed it, so the gap between
R@10 and \code{cap@75} is what reasoning could still gain, and the gap from
\code{cap@75} to 100\% can only be closed by retrieval. Reporting recall without
this ceiling would make a retrieval failure and a reasoning failure look the same.

% =============================================================== 7 results ==
\section{Results}

\subsection{Main comparison}

\begin{table}[H]
\centering
\begin{tabular}{@{}lccc@{}}
\toprule
\textbf{Method} & \textbf{R@1} & \textbf{R@5} & \textbf{R@10}\\
\midrule
CLIP (image-clip)   & \textbf{26.0\%} & 29.0\%          & 33.0\%\\
2 stages            & 23.0\%          & \textbf{46.0\%} & 53.0\%\\
4 stages            & 15.0\%          & 38.0\%          & \textbf{56.0\%}\\
\bottomrule
\end{tabular}
\caption{100 identities out of the 496-identity subset
(Section~\ref{sec:r-dataset}), pool of 75 candidates, zero-shot, one descriptive
phrasing each. At $n=100$ a percentage point is one identity.}
\label{tab:r-main}
\end{table}

Table~\ref{tab:r-main} gives three results.

First, the reasoning layer improves recall at depth. The four-stage funnel returns
the right person in the top 10 for 56 of 100 queries against CLIP's 33, that is,
23 more identities, from reranking alone, with no fine-tuning and no change to
retrieval. At R@5 the gain of the four-stage funnel is smaller, 9 points (38.0\%
against 29.0\%); the two-stage variant gains 17 points there (46.0\% against
29.0\%).

Second, the funnel is worse at the first rank. CLIP is the best method at R@1 (26
against 23 and 15). This is not a rounding effect: the funnel is consistently
better at placing the right person somewhere in the top ten and consistently worse
at placing them first. Section~\ref{sec:r-lost} shows where this happens.

Third, the refine round trades precision for depth. Two stages are better at R@1
and R@5, four stages at R@10. The second round moves into the top 10 candidates
that the coarse round had placed below it, rescuing three more identities, while
costing eight identities at R@1 and eight at R@5. Which of the two rounds should
decide rank 1 remains an open question of the project: the coarse ranker sees 75
candidates with sparse evidence, and the refine ranker sees 25 with dense
evidence.

\subsection{What the funnel changed: rescued and lost}
\label{sec:r-lost}
A recall figure combines two opposite effects. Table~\ref{tab:r-lost} separates
them: of the 33 identities CLIP already had in its top 10, how many survived; and
of the 67 it missed, how many the funnel recovered.

\begin{table}[H]
\centering
\begin{tabular}{@{}lcccc@{}}
\toprule
\textbf{Record} & \textbf{n} & \textbf{Lost} & \textbf{Rescued} & \textbf{No pool}\\
\midrule
2 stages & 100 & 5/33 \ (15\%) & 25/67 \ (37\%) & 3/100\\
4 stages & 100 & 5/33 \ (15\%) & 28/67 \ (42\%) & 3/100\\
\bottomrule
\end{tabular}
\caption{Lost: inside CLIP's top 10, not in the funnel's. Rescued: outside CLIP's
top 10, inside the funnel's. No pool: the ground truth was not among the 75
retrieved crops at all, so no reranking could have returned it.}
\label{tab:r-lost}
\end{table}

The table reproduces Table~\ref{tab:r-main} exactly, which confirms that it
measures the same thing: $33 - 5 + 28 = 56$ for four stages and $33 - 5 + 25 = 53$
for two.

The rescue rate is the main result: 42\% of the identities that CLIP could not
find in its top ten are found by reasoning over attributes, on a pool that CLIP
itself produced. The loss rate is the cost: about one in seven of the identities
already found by retrieval is pushed out. The ceiling is also close: only 3 of 100
ground truths were never retrieved, so 97 of the 100 identities could have been
returned; 56 of them are, and the other 41 (41\% of the sample) are lost somewhere
inside the funnel. The 3 identities outside the pool are not part of these 41,
since no reranking could have returned them. The remaining room for improvement is
therefore in the ranking, not in retrieval.

The loss count is also the same for two and four stages (5 in both), while the
rescue count differs (25 against 28). The refine round adds rescues without adding
losses at the top-10 boundary; its cost appears at R@1 and R@5, inside the top ten.

\subsection{Where the losses happen}
\label{sec:r-lostwhere}
The five lost identities were traced through both rounds. Table~\ref{tab:r-lostwhy}
gives, for each, the rank assigned by CLIP and by each round.

\begin{table}[H]
\centering
\small
\begin{tabular}{@{}ccllp{4.6cm}@{}}
\toprule
\textbf{Candidate} & \textbf{CLIP} & \textbf{Stage 2} & \textbf{Stage 4} &
\textbf{Where it went}\\
\midrule
1 & 1 & 9/73  & 16 & the REFINE round put it 16 of 23\\
2 & 1 & 27/74 & n/a & LOST AT THE COARSE CUT (rank 27 $>$ 25)\\
3 & 1 & 15/75 & 15 & the REFINE round put it 15 of 25\\
4 & 1 & 7/74  & 22 & the REFINE round put it 22 of 24\\
5 & 6 & 30/73 & n/a & LOST AT THE COARSE CUT (rank 30 $>$ 25)\\
\bottomrule
\end{tabular}
\caption{The five identities CLIP had in its top 10 and the funnel did not
return. `Stage 2' is the coarse round's rank out of the candidates it ranked;
`n/a' in stage 4 means the refine round never saw that candidate.}
\label{tab:r-lostwhy}
\end{table}

Two were cut by the coarse round, at ranks 27 and 30 against a cut of 25. These
were near misses, but they cannot be recovered, because stages 3 and 4 never see
them. Three survived the coarse round and were then pushed out by the refine round,
one of them from rank 7 to rank 22. Four of the five were at CLIP rank 1.

That the refine round causes most of the losses is consistent with
Table~\ref{tab:r-main}, where the second round costs R@1 and R@5 while gaining
R@10. It is also why the coarse and refine rounds receive \emph{different}
instructions about absence (Section~\ref{sec:r-coverage}): the coarse cut cannot be
undone and the refine cut can, so the two mistakes are not symmetric and the
prompts should not treat them alike.

\subsection{Is the ranker the bottleneck? A targeted re-run}
\label{sec:r-rerun}
The losses in Table~\ref{tab:r-lostwhy} are ranking failures by construction: the
right person was retrieved, described and ranked, but not high enough. They are
therefore the one set on which the question \emph{``is the small model the
limitation?''} can be asked directly.

The five were re-run with only the ranking stages switched to a larger model
(Gemini Flash 3.8 in place of Flash-Lite 3.5). Everything else was held fixed:
stages 0, 1 and 2A were \emph{replayed} from the original run's saved replies
rather than sent again, so the new ranker read exactly the same evidence as the old
one. Stages 3 and 4A were run again, because their inputs change as soon as the
coarse ranking does.

\begin{table}[H]
\centering
\small
\begin{tabular}{@{}cclll@{}}
\toprule
\textbf{Identity} & \textbf{CLIP} & \textbf{Stage 2} & \textbf{Stage 4} &
\textbf{Result}\\
\midrule
1 & 1 & 9 $\rightarrow$ 7 of 73  & 16 $\rightarrow$ 8 of 24 & FIXED\\
2 & 1 & 27 $\rightarrow$ 1 of 74 & n/a $\rightarrow$ 1 of 25 & FIXED\\
3 & 1 & 15 $\rightarrow$ 3 of 75 & 15 $\rightarrow$ 1 of 25 & FIXED\\
4 & 1 & 7 $\rightarrow$ 14 of 74 & 22 $\rightarrow$ 14 of 25 & STILL WRONG\\
5 & 6 & 30 $\rightarrow$ 1 of 73 & n/a $\rightarrow$ 1 of 24 & FIXED\\
\bottomrule
\end{tabular}
\caption{The five lost identities re-ranked by a larger model at stages 2B and 4B
only. `Stage 2' and `Stage 4' read \emph{before} $\rightarrow$ \emph{after}.}
\label{tab:r-rerun}
\end{table}

Four of the five are recovered, and three of them reach rank 1 of the refine
round. Both coarse-cut losses (candidates 2 and 5) pass the cut easily once a
larger model does the coarse ranking, moving from 27 to 1 and from 30 to 1. The
one that stays wrong, candidate 4, is also the one that the original refine round
demoted the most.

This indicates that these losses came from the capacity of Flash-Lite at the
ranking stages, not from the evidence. The notes were sufficient, but the small
model did not use them well enough.

This result has one important limitation. The five identities were selected using
the ground truth: they are exactly the queries that the original run got wrong
while CLIP got them right. The experiment therefore answers whether the ranker is
the bottleneck \emph{on these failures}. It does not show what the system would
score with the larger model, which only a full sweep can measure, and these four
recoveries must not be added to Table~\ref{tab:r-main}.

\subsection{Case studies}

\paragraph{A lost case (Figure~\ref{fig:r-case-lost}).}
Query: \emph{``A teenage man, with long hair, wearing a short-sleeve red top and
white trousers.''} The ground truth is at CLIP rank 1. The reasoning layer keeps
red tops throughout the ten results it returns, so the garment evidence is used,
but the person it returns first is a different individual in a similar red top,
and the correct identity falls out of the top ten. This is a ranking failure with
correct evidence behind it, of the kind counted in Table~\ref{tab:r-lostwhy}.

\begin{figure}[H]
\centering
\includegraphics[width=\textwidth,height=0.78\textheight,keepaspectratio]{case_lost.png}
\caption{A lost case. Left: three gallery crops of the correct person. Right:
the CLIP ordering and the two funnel rounds, ranks 1 to 5 and 6 to 10.}
\label{fig:r-case-lost}
\end{figure}

\paragraph{The same case after the re-run (Figure~\ref{fig:r-case-lost-fixed}).}
The query and the evidence are identical, and only the two ranking calls are
answered by the larger model. The correct identity returns to the top of both
rounds. Nothing else in the pipeline changed, which makes the comparison direct.

\begin{figure}[H]
\centering
\includegraphics[width=\textwidth,height=0.78\textheight,keepaspectratio]{case_lost_fixed.png}
\caption{The same query after re-ranking with the larger model at stages 2B and
4B only. Stages 0, 1 and 2A were replayed unchanged.}
\label{fig:r-case-lost-fixed}
\end{figure}

\paragraph{A rescued case (Figure~\ref{fig:r-case-rescued}).}
Query: \emph{``A teenage man, with long hair, wearing a short-sleeve purple top
and blue trousers, with a handbag.''} CLIP's top 10 does not contain the right
person, while the funnel's does and places them first. The figure also shows
something that the recall figure cannot: the funnel gathers \emph{several crops of
the same identity} into its ten results without being told that the crops of one
track belong to one person. The grouping follows from the attribute evidence.

\begin{figure}[H]
\centering
\includegraphics[width=\textwidth,height=0.78\textheight,keepaspectratio]{case_rescued.png}
\caption{A rescued case: outside CLIP's top 10, first after reasoning.}
\label{fig:r-case-rescued}
\end{figure}

\subsection{Error analysis}
\label{sec:r-error}
Combining the three tables, the sample of 100 identities divides as follows.

\begin{itemize}
\item 3 were never retrieved. The ground truth is outside the 75-crop pool, so no
      reranking could return it. This can only be fixed in retrieval, and the
      number is small.
\item 56 are returned, 28 of them only because the reasoning layer rescued them.
\item 5 were lost by the funnel after being found by CLIP.
      Section~\ref{sec:r-rerun} shows that four of these five are recovered by a
      larger model reading the \emph{same} evidence, so they are ranking failures,
      not evidence failures.
\item 36 were missed by both methods. They fall into two kinds that can be named
      but have not yet been counted separately at $n=100$: ambiguous attributes,
      where the describer itself is unsure (above all gender on small,
      rear-facing crops), and unread attributes, where a value was never written.
      Accessories are the clearest example: \at{carrying} is unread 28\% of the
      time, more often than its most common value \at{backpack} is read, so a
      query that relies on an accessory favours candidates whose accessory
      \emph{happened} to be described.
\end{itemize}

The failures of the reasoning layer itself, the five lost identities, form the
smallest category and the easiest to address; the large remainder comes from
upstream, from what the index contains.

% =========================================================== 8 ablations ====
\section{Ablations and Design Decisions}
\label{sec:r-ablation}

Almost every mechanism in Section~\ref{sec:r-reasoning} was added in response to a
measured failure rather than planned in advance, and most of them can still be
switched off. This section states, for each mechanism, what it changed and how
strong the evidence is, separating what was \emph{measured} from what is only
\emph{switchable}.

Two limitations apply to all of these numbers.

\paragraph{The sample behind a prompt measurement is a batch, not a sweep.} The
prompt-level numbers below were collected while developing each stage, on batches
of 25 to 75 candidates. They are large, categorical differences (19 to 0, 0 of 82)
rather than narrow ones, which is why they are reported, but they are not recall
measurements over 100 identities.

\paragraph{The model fixes a style per call.} Two batches with the same prompt,
run one after the other, produced 0 of 73 and 73 of 73 notes carrying a boundary
citation. The convention a call adopts at its first contested note holds for the
rest of that call. Any A/B comparison of prompt wording must therefore be read
across several calls, since a single-call comparison of two prompts may only
reflect this style lock.

% Longtable (2026-09-30): the table is taller than a page and a [H] float ran into the bottom margin.
\begin{longtable}{@{}P{4.0cm}P{5.7cm}P{5.4cm}@{}}
\caption{What each mechanism changed, and how strong the evidence is.}\label{tab:r-ablation}\\
\toprule
\textbf{Mechanism} & \textbf{What it changed} & \textbf{Evidence}\\
\midrule
\endfirsthead
\multicolumn{3}{l}{\textit{(continued from Table~\ref{tab:r-ablation})}}\\
\toprule
\textbf{Mechanism} & \textbf{What it changed} & \textbf{Evidence}\\
\midrule
\endhead
\bottomrule
\endlastfoot
Citation rule (2A): never name a mechanism without its figure
& Invented causes disappeared; every named mechanism carries a number.
& \textsc{measured}: 19 invented causes $\rightarrow$ 0; mechanism-without-figure
  0 of 82.\\
\addlinespace
Conditions channel (2A): render the per-region quality figures and pair them with
the value they explain
& The reader began citing capture conditions instead of only counts.
& \textsc{measured}: notes citing a condition with its figure 0 $\rightarrow$ 53
  $\rightarrow$ 73 of 75 across two prompt revisions.\\
\addlinespace
Reliability, not direction (2B)
& Stopped the ranker from opening every row with attributes the whole pool
  shares.
& \textsc{measured}: 69 of 75 rows had opened with the same shared-attribute
  phrase.\\
\addlinespace
Rarity table (2B/4B): state what each query value is worth in the index
& Gave the ranker a rarity channel it did not otherwise have.
& \textsc{measured} as a failure without it: 40 of 40 reasons led with the lower
  colour; gender appeared in 2, while gender split the pool 44/31.\\
\addlinespace
Coverage floor, round-dependent
& Protects a present-but-weak reading from the \emph{unrecoverable} coarse cut
  only.
& \textsc{reasoned}: the coarse cut cannot be undone and the refine cut can.
  Not A/B tested.\\
\addlinespace
Dropping \code{stab} from the decisive gate
& The floor applies to rare values even when the describer writes them
  unsteadily.
& \textsc{measured}: on one query every rare value scored $\text{stab}<0.20$,
  nothing was called decisive, and the floor never applied.\\
\addlinespace
Stage 1/3 restricted to query-constrained attributes
& Stopped the selector from spending slots on attributes the query never
  mentioned.
& \textsc{reasoned} from the selections; not A/B tested.\\
\addlinespace
CLIP score withheld from 2B/4B
& The ranker orders on evidence alone.
& \textsc{measured}: AUC 0.601 within the pool, adjacent-rank gap 2.7\% of one
  sd (Table~\ref{tab:r-clipscore}).\\
\addlinespace
The refine round (4 stages vs 2)
& $+3$ identities at R@10, $-8$ at R@1 and $-8$ at R@5.
& \textsc{measured}, $n=100$ (Table~\ref{tab:r-main}); both records come from the
  same run, so the comparison is paired.\\
\addlinespace
Model at the ranking stages
& 4 of 5 ranking failures recovered on identical evidence.
& \textsc{measured} on 5 oracle-selected queries (Table~\ref{tab:r-rerun}).\\
\end{longtable}

\paragraph{Switchable but not yet measured.} The pipeline has one flag per
experiment so that these comparisons can be run: \code{LAB\_QUALITY\_POPULATION}
(region quartiles over the candidate crops rather than over every frame of their
tracks), \code{LAB\_PAIR\_PROFILE} (the per-pair statistics block),
\code{LAB\_FRAME\_AGREEMENT} (the \code{agrees} line under each quoted frame),
\code{LAB\_TYPICAL\_GUARD}, and \code{VIEW\_MIN\_UNUSUAL} (how unusual a frame must
be before a \code{view} line is printed). They are separate flags rather than one
switch on purpose: the quality-population flag \emph{corrects} an earlier error,
while the others \emph{add} signal, and a result obtained with all of them on
would not show which one caused it.

{\sloppy Three further settings are configuration rather than prompt design and are also
untested: \code{CANDIDATES\_PER\_CALL} (25, the source of the output truncation
discussed in Section~\ref{sec:r-engineering}), \code{RANK\_KEEP\_EXTRA} (how many
rows past the cut the ranker must write), and the thinking level, which is HIGH at
every stage and uses roughly 70\% of the output budget on a 2A call.\par}

\paragraph{The ablation this part could not run.} The most informative experiment,
the whole sweep on the larger model, is limited by quota rather than by design.
Section~\ref{sec:r-rerun} is the part of it that could be afforded.

% ============================================================== 9 tooling ===
\section{Evaluation Tooling and Reproducibility}

Two properties of the measurement setup are reported here because they changed
results.

\paragraph{Records are rebuilt from replies, not trusted.} The sweep writes one
JSON record per configuration and one folder per query holding every stage's
rendered input and reply. Running two notebooks against the same record file
produced a \emph{lost update}: each held the whole file in memory and wrote it
back, so the later save silently dropped the other's queries, while the file still
parsed and could still be scored. The analysis notebook therefore ignores the
record files and rebuilds every ranking from the per-query folders, mapping the
ranker's numeric labels back to crops by aligning the CLIP-similarity sequence in
the saved prompt with the retrieval cache, position by position. When an alignment
check fails, the query is \emph{skipped and named}, never guessed, because
attaching a reply's labels to the wrong candidate list would produce a complete,
plausible and wrong record.

\paragraph{Every rate is reported with its denominator and its ceiling.} Queries
with no ground-truth track are counted and excluded rather than silently dropped,
since a dropped query inflates every rate by shrinking the denominator.
\code{cap@75} is stored per query rather than as one rate, so any later subset can
recompute its own ceiling.

The supporting tools built for this analysis are: a per-identity viewer showing
the ground truth's CLIP rank and its rank after each stage together with the
ranker's reason; a \emph{funnel-lost} report listing the identities that CLIP had
in its top 10 and the funnel did not (the source of Tables~\ref{tab:r-lost}
and~\ref{tab:r-lostwhy}); a coverage table over all evaluated identities; and a
work-queue notebook that classifies every full-pool query not yet evaluated as
\emph{CLIP hit}, \emph{miss but inside the pool} (what reranking can still gain)
or \emph{miss outside the pool} (to be fixed in retrieval).

% ============================================================ 9 limitations =
